\begin{frame}
    \frametitle{Styling Text}
    \framesubtitle{Commands}
    \begin{center}
    % \begin{tabular}{l l}
    %     \textbf{Command}
    %     & \textbf{Example}\\
    % \toprule[1.2pt]
    %     \lstinline|\\emph\{\}|          
    %     & \emph{Emphasized}\\
    % \midrule%
    %     \lstinline|\\textbf\{\}|          
    %     & \textbf{Bold}\\
    % \midrule%
    %     \lstinline|\\textit\{\}|       
    %     & \textit{Italic}\\
    % \midrule
    %     \lstinline|\\textsc\{\}|       
    %     & \textsc{SmallCaps}\\
    % \midrule
    %     \lstinline|\\textsl\{\}|       
    %     & \textsl{Slanted}\\
    % \midrule
    %     \lstinline|\\textsf\{\}|       
    %     & \textsf{Sans-Serif}\\
    % \midrule
    %     \lstinline|\\textrm\{\}|       
    %     & \textrm{Roman}\\
    % \midrule
    %     \lstinline|\\texttt\{\}|       
    %     & \texttt{Typewriter}\\
    % \midrule
    %     \lstinline|\\underline\{\}|       
    %     & \underline{Underlined}\\
    % \end{tabular}
    \onslide<2->%
    \begin{tabular}{@{}l l@{}}
        \textbf{Command}
        & \textbf{Example}\\
    \toprule[1.2pt]
        \lstinline|\\emph\{\}|          
        & \emph{Emphasized}\\
    \midrule%
        \lstinline|\\textit\{\}|       
        & \textit{Italic}\\
    \midrule
        \lstinline|\\textsl\{\}|       
        & \textsl{Slanted}\\
    \end{tabular}%
    \onslide<3->%
    \quad%
    \begin{tabular}{@{}l l@{}}
        \textbf{Command}
        & \textbf{Example}\\
    \toprule[1.2pt]
        \lstinline|\\textbf\{\}|          
        & \textbf{Bold}\\
    \midrule
        \lstinline|\\textsc\{\}|       
        & \textsc{SmallCaps}\\
    \midrule
        \lstinline|\\textsf\{\}|       
        & \textsf{Sans-Serif}\\
    \end{tabular}%
    \onslide<4->%
    \quad%
    \begin{tabular}{@{}l l@{}}
        \textbf{Command}
        & \textbf{Example}\\
    \toprule[1.2pt]
        \lstinline|\\textrm\{\}|       
        & \textrm{Roman}\\
    \midrule
        \lstinline|\\texttt\{\}|       
        & \texttt{Typewriter}\\
    \midrule
        \lstinline|\\underline\{\}|       
        & \underline{Underlined}\\
    \end{tabular}
    
    \end{center}
\end{frame}

\begin{frame}
    \frametitle{Styling Text}
    \begin{columns}
        \begin{column}{.59\textwidth}
            \alt<1>{\tcbexinputlst{listings/sections.tex}}%
            {\tcbexinputlst{listings/text-styling.tex}}
        \end{column}
        \begin{column}{.4\textwidth}
            \centering
            \includegraphics<3->[%
                width=0.95\textwidth,
                frame,
                clip=true,
                trim=3cm 23cm 3cm 4cm
            ]{listings/text-styling.pdf}   
        \end{column}
    \end{columns}
\end{frame}

% \begin{frame}
%     \frametitle{Styling Text}
%     \framesubtitle{Remarks}

%     \begin{itemize}
%         \item<+-> Usually only \lstinline|\\emph\{\}| should be used!
%     \end{itemize}

%     \uncover<2->{%
%     \begin{center}
%         \Large
%         Why?
%     \end{center}
%     }

%     \begin{itemize}
%         \item<3-> Emphasized text is always different than surrounding text.
%     \end{itemize}

%     \uncover<4->{%
%     \begin{exblock}
%     \begin{columns}
%         \begin{column}{.68\textwidth}
%             \lstinline|\\emph\{Emphasized text is \\emph\{toggled\}.\}|
%         \end{column}
%         \begin{column}{.3\textwidth}
%             \emph{Emphasized text is \emph{toggled}.}
%         \end{column}
%     \end{columns}
%     \end{exblock}}
% \end{frame}